\centerline{\bf \Huge Виет настигает}

\begin{flushright}
        \scriptsize{}
\end{flushright}

\problem{1}
Используя теорему Виета, угадайте корни квадратных уравнений:

(a) $x^2-5 x+6=0$;

(b) $2 x^2-5 x+3=0$;

(c) $x^2-(2 a+4) x+a^2+4 a=0$.

\problem{2}
Не вычисляя корней уравнения $3 x^2+4 x-1=0$, найдите

(a) $x_1^2+x_2^2$,

(b) $x_2^3 x_1+x_1^3 x_2$,

(c) $x_2^3+x_1^3$.

\problem{3}
Пусть $x_1$ и $x_2$ - корни уравнения $2 x^2-7 x-3=0$. Составьте квадратное уравнение, корнями которого будут являться числа

(a) $x_1-1$ и $x_2-1$;

(b) $\frac{1}{x_1}$ и $\frac{1}{x_2}$;

(c) $x_1 x_2^2$ и $x_2 x_1^2$;

(d) $\frac{x_1}{x_2}+1$ и $\frac{x_2}{x_1}+1$.

\problem{4}
Известно, что корни уравнения $x^2+p x+q=0$ - целые числа, а $p$ и $q$ - простые числа. Найдите $p$ и $q$.

\problem{5}
Квадратный трёхчлен $f(x)=a x^2+b x+c$ принимает в точках $1 / a$ и $c$ значения разных знаков. Докажите, что корни трёхчлена $f(x)$ имеют разные знаки.
